% Part: first-order-logic
% Chapter: axiomatic-deduction
% Section: provability-quantifiers

% verification of properties of provability needed for maximally
% consistent sets in the completeness chapter.

\documentclass[../../../include/open-logic-section]{subfiles}

\begin{document}

\olfileid{fol}{axd}{qpr}

\olsection{\usetoken{S}{derivability} અને પરિમાણકો}

\begin{explain}
  પૂર્ણતાના પ્રમેય માટે એ પણ જરૂરી છે કે સ્વયંસિદ્ધિમૂલક નિષ્પત્તિઓ
  આ વિભાગમાં સ્થાપેલાં~$\Proves$ વિશેનાં તથ્યો આપે.
\end{explain}

\begin{thm}
\ollabel{thm:strong-generalization} જો $c$ એ~$\Gamma$ કે $!A(x)$માં ન
આવતું !!a{constant} હોય અને $\Gamma \Proves !A(c)$, તો
$\Gamma \Proves \lforall[x][!A(x)]$.
\end{thm}

\begin{proof}
નિગમન પ્રમેય વડે $\Gamma \Proves \ltrue \lif !A(c)$. કારણ કે $c$ એ
$\Gamma$ કે~$\ltrue$માં આવતું નથી, તેથી
$\Gamma \Proves \ltrue \lif \lforall[x][!A(x)]$ મળે છે. હવે
$\ltrue$ સ્વયંસિદ્ધિ હોવાથી મોડસ પોનેન્સ વડે
$\Gamma \Proves \lforall[x][!A(x)]$.
\end{proof}

\sourcecorrection{OLAX-011}{મૂળની
\texttt{provability-quantifiers.tex}, પંક્તિ~29માં અગાઉ અને પછી વપરાતા
ભાષા-ચિહ્ન~$\ltrue$ને બદલે કાચું \texttt{\string\top} લખાયું છે. અહીં
એક જ સ્વયંસિદ્ધિ-ચિહ્ન~$\ltrue$ રાખ્યું છે.}

\sourcecorrection{OLAX-012}{મૂળની
\texttt{provability-quantifiers.tex}, પંક્તિઓ 30--31માં
$\Gamma\Proves\ltrue\lif\lforall[x][!A(x)]$માંથી સર્વત્રીકૃત પરિણામને
``નિગમન પ્રમેય ફરી'' વડે મેળવ્યું છે. જરૂરી પગલું એ છે કે $\ltrue$
સ્વયંસિદ્ધિ હોવાથી તેને અને પ્રેરણને મોડસ પોનેન્સમાં વાપરવું; અહીં એ
પ્રમાણન લખ્યું છે.}

\begin{prop}
\ollabel{prop:provability-quantifiers}
\begin{tagenumerate}{prvEx,prvAll}
\tagitem{prvEx}{$!A(t) \Proves \lexists[x][!A(x)]$.}{}

\tagitem{prvAll}{$\lforall[x][!A(x)] \Proves !A(t)$.}{}
\end{tagenumerate}
\end{prop}

\begin{proof}
\begin{tagenumerate}{prvEx,prvAll}
\tagitem{prvEx}{\olref[qua]{ax:q2} અને નિગમન પ્રમેય વડે.}{}
\tagitem{prvAll}{\olref[qua]{ax:q1} અને નિગમન પ્રમેય વડે.}{}
\end{tagenumerate}
\end{proof}

\end{document}
